\documentclass[12pt, a4paper]{article}

% ── Packages ──────────────────────────────────────────────────────────────────
\usepackage[utf8]{inputenc}
\usepackage[T1]{fontenc}
\usepackage[margin=2.5cm]{geometry}
\usepackage{amsmath, amssymb, amsthm}
\usepackage{graphicx}
\usepackage{booktabs}
\usepackage{hyperref}
\usepackage{cleveref}
\usepackage{algorithm}
\usepackage{algpseudocode}
\usepackage{xcolor}
\usepackage{tikz}
\usetikzlibrary{shapes.geometric, arrows.meta, positioning, fit, backgrounds}
\usepackage{subcaption}
\usepackage{enumitem}
\usepackage{multirow}
\usepackage{tabularx}
\usepackage{lmodern}
\usepackage{microtype}
\usepackage{setspace}
\usepackage[numbers, sort&compress]{natbib}

\onehalfspacing

% ── Macros ────────────────────────────────────────────────────────────────────
\newcommand{\piA}{\pi_A}
\newcommand{\piB}{\pi_B}
\newcommand{\piC}{\pi_C}
\newcommand{\piM}{\pi_M}
\newcommand{\sstar}{s^{*}}
\newcommand{\Enc}{E}
\newcommand{\goalvec}{\mathbf{g}}
\newcommand{\RR}{\mathbb{R}}
\newcommand{\calL}{\mathcal{L}}
\newcommand{\calD}{\mathcal{D}}
\newcommand{\calG}{\mathcal{G}}
\newcommand{\calS}{\mathcal{S}}
\newcommand{\calA}{\mathcal{A}}
\newcommand{\arm}[1]{\textsc{Arm}_{#1}}
\newtheorem{definition}{Definition}
\newtheorem{remark}{Remark}

% ── Colours ───────────────────────────────────────────────────────────────────
\definecolor{aliceblue}{RGB}{0,114,189}
\definecolor{bobred}{RGB}{217,83,25}
\definecolor{charliegreen}{RGB}{32,134,48}
\definecolor{managerviolet}{RGB}{126,47,142}

% ── Hyperref Setup ────────────────────────────────────────────────────────────
\hypersetup{
    colorlinks=true,
    linkcolor=blue!60!black,
    citecolor=green!50!black,
    urlcolor=blue,
    pdftitle={Dual-Arm Robotic Manipulation via Multi-Agent Asymmetric Self-Play and HRL},
    pdfauthor={},
}

% ══════════════════════════════════════════════════════════════════════════════
\begin{document}

\title{%
  \textbf{Dual-Arm Robotic Manipulation via\\
  Multi-Agent Asymmetric Self-Play and\\
  Hierarchical Reinforcement Learning}\\[0.5em]
  \large Implementation Plan \& Technical Design Document
}
\author{}
\date{\today}
\maketitle
\tableofcontents
\newpage

% ══════════════════════════════════════════════════════════════════════════════
\section{Introduction}
\label{sec:intro}
% ══════════════════════════════════════════════════════════════════════════════

Dual-arm robotic manipulation represents one of the most challenging open problems in robot
learning. Tasks such as assembly, bi-manual grasping, and dexterous object re-arrangement
require tight spatiotemporal coordination between two arms that share a workspace, yet each
arm must also be competent in isolation. Training such a system from scratch with sparse
external reward is notoriously difficult due to the exponential state-space explosion,
compounded credit-assignment problems, and the need to discover role specialisation without
explicit supervision.

This document presents a full design and implementation plan for a framework that addresses
these challenges by combining three complementary ideas from recent literature:

\begin{enumerate}[label=(\roman*)]
    \item \textbf{Asymmetric Self-Play (ASP)} \cite{sukhbaatar2018intrinsic,plappert2021asymmetric}
          for unsupervised pre-training of individual arm sub-policies without any external reward.
    \item \textbf{Hierarchical goal-embedding via self-play} \cite{sukhbaatar2018learning} to learn
          a compact, reusable latent goal space that bridges arm-level and manager-level control.
    \item \textbf{Manager-level meta-ASP} to train a high-level coordination policy that issues
          joint goal pairs to both frozen arms, again without external reward.
\end{enumerate}

The simulator is NVIDIA IsaacLab \cite{isaaclab2023}, built on PhysX, which is leveraged
as a physics oracle to replace hand-crafted penalty formulas with exact collision and
proximity queries.

% ══════════════════════════════════════════════════════════════════════════════
\section{Theoretical Background}
\label{sec:background}
% ══════════════════════════════════════════════════════════════════════════════

\subsection{Markov Decision Processes and Goal-Conditioned RL}
\label{subsec:mdp}

We model the environment as a goal-augmented Markov Decision Process
$\mathcal{M} = \langle \calS, \calA, \mathcal{P}, \mathcal{R}, \calG \rangle$
where $\calS$ is the state space, $\calA$ the action space,
$\mathcal{P}: \calS \times \calA \times \calS \mapsto [0,1]$ the transition probability,
$\calG \subseteq \calS$ the goal space, and
$\mathcal{R}: \calS \times \calG \mapsto \RR$ the goal-conditioned reward
\cite{kaelbling1993learning,plappert2021asymmetric}.
A goal-conditioned policy $\pi(a \mid s, g)$ maps the current state and a goal
to a distribution over actions. A goal $g$ is considered achieved when the state
$s_t$ is within a success threshold $\epsilon$ of $g$ under distance $D$:
\begin{equation}
    D(s_t, g) \leq \epsilon.
\end{equation}

\subsection{Asymmetric Self-Play (ASP)}
\label{subsec:asp}

Asymmetric Self-Play, introduced by \citet{sukhbaatar2018intrinsic}, is an
\emph{unsupervised} curriculum generation scheme in which two agents — Alice and Bob —
play an adversarial game within a shared environment.

\paragraph{Game structure.}
Alice acts for $T_A$ steps from initial state $s_0$, reaching state $\sstar$, which becomes
Bob's goal. Bob then starts from $s_0$ and must reach $\sstar$ within $T_B \geq T_A$ steps.
The reward structure is
\begin{align}
    R_B &= -\gamma\, t_B, \label{eq:rb_og}\\
    R_A &= \gamma \max(0,\, t_B - t_A), \label{eq:ra_og}
\end{align}
where $t_A$, $t_B$ are the steps each agent requires and $\gamma$ is a scaling coefficient.
This reward incentivises Alice to propose goals that are just beyond Bob's current ability,
automatically constructing a curriculum without any human specification
\cite{sukhbaatar2018intrinsic}.

\paragraph{OpenAI Robotic Manipulation variant.}
\citet{plappert2021asymmetric} extend ASP to a table-top manipulation setting with a robot
arm. Alice physically rearranges objects to propose a goal workspace configuration; Bob must
replicate it. The multi-goal game runs up to 5 goals per episode, allowing Bob to build
environmental context across goals. A key addition is \emph{Alice Behavioral Cloning} (ABC):
since Alice's trajectory always constitutes a valid solution to Bob's goal, Bob is jointly
trained with a PPO-style clipped BC loss on Alice's demonstrations:
\begin{equation}
    \calL_{\mathrm{ABC}} = -\mathbb{E}_{(s_t,a_t,g_t)\in\calD_{\mathrm{BC}}}
    \left[
        \mathrm{clip}\!\left(
            \frac{\pi_B(a_t \mid s_t, g_t; \theta)}{\pi_B(a_t \mid s_t, g_t; \theta_{\mathrm{old}})},
            1-\epsilon_{\mathrm{clip}},\, 1+\epsilon_{\mathrm{clip}}
        \right)
    \right], \label{eq:abc}
\end{equation}
where $\hat{A}=1$ is used as the advantage (treating every Alice demonstration as a positive
example). Only demonstrations where Bob \emph{failed} are retained, discarding sub-optimal
Alice trajectories for trivially solved goals \cite{plappert2021asymmetric}.

\subsection{Goal Embeddings and Hierarchical Self-Play (HSP)}
\label{subsec:hsp}

\citet{sukhbaatar2018learning} introduce a \emph{hierarchical} extension of ASP (HSP) in
which a goal encoder $\Enc$ maps state pairs to a low-dimensional embedding
$\goalvec_t = \Enc(\sstar, s_t^B) \in \RR^K$.
Bob's policy is then conditioned on this embedding:
\begin{equation}
    a_t^B = \pi_B'(s_t^B,\, \goalvec_t) = \pi_B'(s_t^B,\, \Enc(\sstar, s_t^B)). \label{eq:bob_enc}
\end{equation}
The encoder is trained jointly with Alice and Bob during self-play. After pre-training, a
high-level \emph{Charlie} policy $\pi_C$ is introduced:
\begin{equation}
    \goalvec_t = \pi_C(s_t), \label{eq:charlie}
\end{equation}
which operates at a coarser timescale $T_C \approx T_A$, issuing goal vectors every $T_C$
environment steps. Charlie is then trained on the target task using external reward while
Bob and $\Enc$ are fine-tuned. The HSP paper further notes, in its discussion, that a natural
extension is a \emph{meta-Alice / meta-Bob} pair playing at the Charlie level, passing
sub-goals to lower-level Alice and Bob — precisely the structure proposed in this work.

\subsection{Proximal Policy Optimisation (PPO)}
\label{subsec:ppo}

All policies in this framework are trained with PPO \cite{schulman2017proximal}. The
clipped surrogate objective is:
\begin{equation}
    \calL_{\mathrm{PPO}} = \mathbb{E}_t\!\left[
        \min\!\left(
            r_t(\theta)\hat{A}_t,\;
            \mathrm{clip}(r_t(\theta), 1-\epsilon, 1+\epsilon)\hat{A}_t
        \right)
    \right], \label{eq:ppo}
\end{equation}
where $r_t(\theta) = \pi_\theta(a_t \mid s_t) / \pi_{\theta_{\mathrm{old}}}(a_t \mid s_t)$
and $\hat{A}_t$ is the Generalised Advantage Estimate (GAE).

\subsection{Hierarchical Reinforcement Learning}
\label{subsec:hrl}

Hierarchical RL decomposes a long-horizon task into a manager that sets sub-goals and
one or more workers that achieve them \cite{vezhnevets2017feudal,nachum2018data}.
In Feudal Networks \cite{vezhnevets2017feudal}, the manager operates at a coarser timescale
and communicates with workers via a goal vector in a learned embedding space —
structurally identical to the Charlie architecture in \citet{sukhbaatar2018learning}.

% ══════════════════════════════════════════════════════════════════════════════
\section{Proposed Framework}
\label{sec:framework}
% ══════════════════════════════════════════════════════════════════════════════

\subsection{Overview}
\label{subsec:overview}

The framework operates in two sequential training phases with a frozen deployment stage:

\begin{description}
    \item[Phase 1 — Single-Arm Pre-training via ASP.]
          A shared sub-policy $\piB'$ and goal encoder $\Enc$ are trained via single-arm
          ASP. Both arms share identical parameters.

    \item[Phase 2 — Manager Training via Meta-ASP.]
          The sub-policies are frozen. A manager $\piM$ is trained via a second ASP layer
          (meta-Alice / meta-Bob) that operates on joint workspace goals. No external reward
          is used.

    \item[Deployment.]
          The manager issues synchronised goal pairs $[\goalvec_1, \goalvec_2]$ to both
          frozen arms every $T_C$ steps.
\end{description}

\subsection{Notation}
\label{subsec:notation}

\begin{table}[h!]
\centering
\caption{Key notation used throughout this document.}
\begin{tabular}{ll}
\toprule
Symbol & Meaning \\
\midrule
$s_t^{(i)}$ & State observation of arm $i \in \{1,2\}$ at time $t$ \\
$\goalvec^{(i)} \in \RR^K$ & Goal embedding for arm $i$ \\
$\Enc$ & Shared goal encoder, maps $(s^*, s_t) \mapsto \RR^K$ \\
$\piB'$ & Shared frozen arm sub-policy, $a_t = \piB'(s_t, \goalvec_t, r_{\mathrm{id}})$ \\
$\piM$ & Manager policy, outputs $[\goalvec^{(1)}, \goalvec^{(2)}] \in \RR^{2K}$ \\
$\piA^{(1)}$ & meta-Alice: manager-level goal proposer \\
$\piA^{(2)}$ & meta-Bob: the manager $\piM$ in goal-solving mode \\
$T_A$, $T_B$ & Alice and Bob step budgets in Phase~1 \\
$T_C$ & Manager goal-update period ($T_C \approx T_A$) \\
$K$ & Goal embedding dimension \\
$r_{\mathrm{id}} \in \{0,1\}$ & Role identifier appended to arm observation \\
\bottomrule
\end{tabular}
\label{tab:notation}
\end{table}

% ══════════════════════════════════════════════════════════════════════════════
\section{Phase 1: Single-Arm Pre-training via ASP}
\label{sec:phase1}
% ══════════════════════════════════════════════════════════════════════════════

\subsection{Observation Structure}
\label{subsec:obs_phase1}

To enable a shared policy to serve as both $\arm{1}$ and $\arm{2}$ at deployment without
confusion about which joints it controls, the observation is structured with a fixed global
ordering and a \emph{role identifier}:

\begin{equation}
    o_t^{(i)} = \bigl[\underbrace{s_t^{(i)}}_{\text{own joints}},\;
                       \underbrace{s_t^{(j)}}_{\text{other joints (padded)}},\;
                       \goalvec_t^{(i)},\;
                       r_{\mathrm{id}}^{(i)}\bigr],
    \quad j \neq i. \label{eq:obs}
\end{equation}

During Phase~1, only one arm is present. The ``other joints'' slot is \emph{zero-padded}.
To prevent the policy from trivially ignoring the role identifier (which would be constant
if always set to 0), each pre-training episode randomly samples $r_{\mathrm{id}} \in \{0,1\}$
and correspondingly places the arm's real joint values in the matching position:

\begin{align}
    r_{\mathrm{id}} = 0 \;\Rightarrow\; o_t &= [\,s_t^{\text{arm}},\; \mathbf{0},\; \goalvec_t,\; 0\,], \label{eq:rid0}\\
    r_{\mathrm{id}} = 1 \;\Rightarrow\; o_t &= [\,\mathbf{0},\; s_t^{\text{arm}},\; \goalvec_t,\; 1\,]. \label{eq:rid1}
\end{align}

The simulator always routes the policy's action output to the single physical arm, regardless
of $r_{\mathrm{id}}$. The policy must learn to use $r_{\mathrm{id}}$ to identify which block
of the observation corresponds to its own proprioception; otherwise it fails half the time
(conditioning on zeros). At deployment the zero-padding is replaced by the second arm's real
state, providing natural collision awareness.

\subsection{Alice and Bob Sub-policies}
\label{subsec:phase1_policies}

Alice and Bob share the same network architecture but have separate parameters
$\theta_A$ and $\theta_B$. Bob's policy additionally passes through the goal encoder:

\begin{equation}
    a_t^B = \piB'(s_t^B,\; \Enc(\sstar, s_t^B),\; r_{\mathrm{id}}).
\end{equation}

Alice's policy is $\piA(s_t^A, s_0)$ following \citet{sukhbaatar2018intrinsic} and
\citet{sukhbaatar2018learning}. Alice is given entropy regularisation to encourage
diverse task proposals:
\begin{equation}
    \calL_A = \mathbb{E}\bigl[-R_A - \beta\, \mathcal{H}(\piA(s_t^A))\bigr]. \label{eq:alice_loss}
\end{equation}

\subsection{Multi-Goal Game Structure}
\label{subsec:multi_goal}

Following \citet{plappert2021asymmetric}, each episode runs up to $N_G = 5$ self-play goals
sequentially. Bob's last reached state becomes the initial state for the next goal, allowing
exploration of states far from $s_0$:

\begin{enumerate}[label=\arabic*.]
    \item Alice acts for $T_A$ steps; her final state $\sstar$ becomes Bob's goal.
    \item Bob attempts to reach $\sstar$ from the current state within $T_B$ steps.
    \item If Bob succeeds, the episode continues with the next goal; otherwise it terminates.
    \item After $N_G$ goals (or early termination), rewards are assigned and policies updated.
\end{enumerate}

\subsection{Alice Behavioral Cloning at the Arm Level}
\label{subsec:abc_arm}

Bob's full training objective combines the PPO loss with the ABC loss from
\citet{plappert2021asymmetric}:
\begin{equation}
    \calL_B = \calL_{\mathrm{PPO}}^B + \beta\, \calL_{\mathrm{ABC}}, \label{eq:bob_loss}
\end{equation}
where $\calL_{\mathrm{ABC}}$ is defined in \cref{eq:abc}.
Only Alice trajectories where Bob \emph{failed} are added to the BC replay buffer
$\calD_{\mathrm{BC}}$, discarding suboptimal demonstrations for already-mastered goals.

\subsection{Goal Encoder Architecture}
\label{subsec:encoder}

Two encoder variants are considered, following \citet{sukhbaatar2018learning}:

\begin{description}
    \item[Difference encoder:]
          $\Enc(\sstar, s_t^B) = \phi(\sstar) - \phi(s_t^B)$,
          where $\phi$ is a two-layer MLP without final non-linearity.
    \item[Absolute encoder:]
          $\Enc(\sstar, s_t^B) = \phi(\sstar)$,
          which ignores the current state.
\end{description}

The goal embedding dimension $K$ must be large enough to retain geometric information about
end-effector targets (recommended $K \in [8, 16]$, larger than the $K=2$ used for 2D Ant
tasks in \citet{sukhbaatar2018learning}).

\subsection{Phase 1 Network Architecture}
\label{subsec:network_phase1}

Following the OpenAI manipulation architecture \cite{plappert2021asymmetric}:

\begin{itemize}
    \item \textbf{Robot state embedding:} Linear(256) + LayerNorm on joint positions and
          gripper positions separately.
    \item \textbf{Object state embedding:} Permutation-invariant embedding (concatenate per-object
          features, embed, max-pool over objects) with 512 units.
    \item \textbf{Goal state embedding:} Linear(256) + LayerNorm on $\goalvec$.
    \item \textbf{Trunk:} Sum of embeddings $\to$ ReLU $\to$ MLP $\to$ LSTM.
    \item \textbf{Output heads:} Separate policy head (action distribution) and value head.
\end{itemize}

\subsection{Phase 1 Termination and Freezing}
\label{subsec:phase1_stop}

Phase~1 is considered complete when Alice's proposed task distribution stabilises and Bob
achieves high success rates across a diverse held-out set of single-arm manipulation tasks
(push, flip, pick-and-place). At this point:

\begin{itemize}[nosep]
    \item $\piB'(\cdot;\, \theta_B)$ is \textbf{frozen}. Parameters are never updated again.
    \item $\Enc(\cdot;\, \theta_E)$ is \textbf{frozen}.
    \item $\piA$ is discarded (it was only a curriculum generator).
\end{itemize}

% ══════════════════════════════════════════════════════════════════════════════
\section{Phase 2: Manager Training via Meta-ASP}
\label{sec:phase2}
% ══════════════════════════════════════════════════════════════════════════════

\subsection{Manager Architecture}
\label{subsec:manager_arch}

The manager $\piM$ maps the full bimanual state and an optional task context to a joint
goal pair:

\begin{equation}
    [\goalvec^{(1)}, \goalvec^{(2)}] = \piM\!\left([s_t^{(1)},\; s_t^{(2)},\; z_{\mathrm{task}}]\right)
    \in \RR^{2K}, \label{eq:manager_output}
\end{equation}

where $z_{\mathrm{task}}$ is an optional task-context vector (empty during meta-ASP
pre-training). The manager uses a \textbf{shared trunk} with two separate output heads:

\begin{align}
    h &= \mathrm{MLP}_{\mathrm{trunk}}\!\left([s_t^{(1)},\, s_t^{(2)},\, z_{\mathrm{task}}]\right),\\
    \goalvec^{(1)} &= \mathrm{Head}_1(h), \quad
    \goalvec^{(2)} = \mathrm{Head}_2(h).
\end{align}

The shared trunk is essential: it forces the manager to jointly reason about both arms
before committing to either goal, preventing the degenerate case where each head ignores
the other arm's state.

The manager issues $[\goalvec^{(1)}, \goalvec^{(2)}]$ every $T_C$ steps simultaneously
to both arms. Both arms execute their frozen sub-policies for $T_C$ steps with fixed goals,
then the manager updates. We set $T_C = T_A$ following \citet{sukhbaatar2018learning},
since Bob is trained to be reliable for goals achievable within $T_A$ steps.

\subsection{Meta-ASP Game Structure}
\label{subsec:meta_asp}

The meta-ASP layer mirrors the arm-level ASP one abstraction level higher. Both
meta-Alice and meta-Bob operate in the same action space: sequences of
$[\goalvec^{(1)}, \goalvec^{(2)}]$ vectors issued to the frozen arm sub-policies.

\begin{description}
    \item[meta-Alice ($\piA^M$):] A manager-like policy that issues joint goal sequences
          to physically rearrange the workspace into a target configuration $\sstar_{\mathrm{ws}}$
          (a full object layout achievable by two arms acting together).

    \item[meta-Bob ($\piM$):] The manager being trained. Starting from the same initial
          state $s_0$, it must issue joint goal sequences to drive the workspace from
          $s_0$ to $\sstar_{\mathrm{ws}}$.
\end{description}

Since meta-Alice operates through the same frozen arms, her trajectory of goal-vector
sequences constitutes a \textbf{valid demonstration} for meta-Bob — the same guarantee
that enables ABC at the arm level \cite{plappert2021asymmetric}.

\paragraph{Reward structure.}
\begin{align}
    R_{\text{meta-Bob}} &= +1 \;\text{if workspace reaches } \sstar_{\mathrm{ws}} \text{ within } D_{\mathrm{ws}},\;
                            0 \;\text{otherwise,} \label{eq:meta_bob_rew}\\
    R_{\text{meta-Alice}} &= 1 - R_{\text{meta-Bob}}, \label{eq:meta_alice_rew}
\end{align}
where the workspace distance $D_{\mathrm{ws}}$ is the sum of per-object Euclidean
(position) and Euler-angle (rotation) distances, identical to the metric used at
the object level in \citet{plappert2021asymmetric} (success thresholds: 0.04\,m,
0.2\,rad).

\paragraph{Goal validation.}
meta-Alice's proposed workspace configuration $\sstar_{\mathrm{ws}}$ is validated before
Bob's turn begins, following \citet{plappert2021asymmetric}~Appendix~A.1:
\begin{enumerate}[label=\arabic*., nosep]
    \item At least one object must have moved relative to $s_0$.
    \item All objects must remain on the workspace surface.
    \item All objects must lie within the reachable placement area of both arms.
\end{enumerate}
Invalid goals trigger an episode reset without reward.

\subsection{ABC at the Manager Level}
\label{subsec:abc_manager}

meta-Alice's goal-vector trajectory $\tau_A^M = \{[\goalvec_0^{(1)}, \goalvec_0^{(2)}], \ldots\}$
is a valid demonstration of how the manager should coordinate the two arms to reach
$\sstar_{\mathrm{ws}}$. The manager's full training objective is:
\begin{equation}
    \calL_M = \calL_{\mathrm{PPO}}^M + \beta_M\, \calL_{\mathrm{ABC}}^M
              + \lambda\, \calL_{\mathrm{collision}} + \mu\, \calL_{\mathrm{conflict}},
    \label{eq:manager_loss}
\end{equation}
where $\calL_{\mathrm{ABC}}^M$ clones meta-Alice's goal-vector sequences (with the same
PPO-style clipping and demonstration filtering as the arm-level ABC), and the last two
terms are simulation-based physical penalties described in \cref{sec:penalties}.

\subsection{Simulation-Based Physical Penalties}
\label{sec:penalties}

Rather than hand-crafting geometric approximations, both penalties are computed directly
from IsaacLab's physics engine, which provides exact collision and proximity information.

\paragraph{Arm-Arm Collision Penalty ($\calL_{\mathrm{collision}}$).}
A \texttt{ContactSensor} is configured to monitor contacts between the link bodies of
$\arm{1}$ and $\arm{2}$. Contact forces are reported by PhysX at every simulation step.
The penalty is accumulated over the entire $T_C$ execution window and normalised:
\begin{equation}
    \calL_{\mathrm{collision}} = \frac{1}{T_C} \sum_{t=1}^{T_C}
        \max_j \|\mathbf{f}_t^{(j)}\|_2,
    \label{eq:collision_penalty}
\end{equation}
where $\mathbf{f}_t^{(j)}$ is the net contact force on link $j$ of $\arm{1}$ due to
contact with $\arm{2}$. This is exact, handles all geometries (including non-convex
link decompositions), and requires no decoder or forward-kinematics approximation.

In IsaacLab:
\begin{verbatim}
from omni.isaac.lab.sensors import ContactSensor, ContactSensorCfg

contact_sensor = ContactSensor(ContactSensorCfg(
    prim_path="{ENV_REGEX_NS}/arm1/.*",
    filter_prim_paths_expr=["{ENV_REGEX_NS}/arm2/.*"],
    track_air_time=False,
))
# In reward computation:
forces = contact_sensor.data.net_forces_w   # [N_envs, N_bodies, 3]
L_col  = forces.norm(dim=-1).max(dim=-1).values  # [N_envs]
\end{verbatim}

\paragraph{Object Conflict Penalty ($\calL_{\mathrm{conflict}}$).}
An object conflict occurs when both arms' end-effectors are simultaneously in close
proximity to the same object. This is detected directly from body position queries:
\begin{equation}
    \calL_{\mathrm{conflict}} = \mathbb{1}\!\left[
        \exists\, k : D(\mathrm{ee}_1, o_k) < d_{\mathrm{prox}}
        \;\wedge\;
        D(\mathrm{ee}_2, o_k) < d_{\mathrm{prox}}
    \right],
    \label{eq:conflict_penalty}
\end{equation}
where $o_k$ is the position of object $k$ and $d_{\mathrm{prox}} \approx 0.08$\,m.
In IsaacLab:
\begin{verbatim}
ee1 = arm1.data.body_pos_w[:, ee1_idx, :]        # [N, 3]
ee2 = arm2.data.body_pos_w[:, ee2_idx, :]        # [N, 3]
obj = objects.data.root_pos_w                     # [N, M, 3]

near1 = (obj - ee1.unsqueeze(1)).norm(dim=-1) < prox_thresh
near2 = (obj - ee2.unsqueeze(1)).norm(dim=-1) < prox_thresh
L_conflict = (near1 & near2).any(dim=-1).float() # [N]
\end{verbatim}

\paragraph{Penalty scheduling.}
The collision weight $\lambda$ is annealed from high to low during manager training to
prevent early degenerate collisions while allowing close cooperation later:
\begin{equation}
    \lambda(t) = \lambda_{\max} \cdot \exp\!\left(-\frac{t}{\tau_\lambda}\right) + \lambda_{\min}.
\end{equation}

% ══════════════════════════════════════════════════════════════════════════════
\section{Algorithms}
\label{sec:algorithms}
% ══════════════════════════════════════════════════════════════════════════════

\subsection{Phase 1: Single-Arm ASP Training}
\label{subsec:alg_phase1}

\begin{algorithm}[H]
\caption{Phase 1 — Single-Arm ASP with Goal Encoder}
\label{alg:phase1}
\begin{algorithmic}[1]
\Require $\theta_A, \theta_B, \theta_E$ \Comment{Initial parameters: Alice, Bob, Encoder}
\Require $\eta$, $\beta$, $N_G$, $T_A$, $T_B$, $T_C$ \Comment{Hyperparameters}
\For{training step $= 1, 2, \ldots$}
    \State $\theta_{A}^{\mathrm{old}} \leftarrow \theta_A$;\quad
           $\theta_{B}^{\mathrm{old}} \leftarrow \theta_B$
    \For{each parallel environment worker}
        \State $r_{\mathrm{id}} \sim \mathrm{Uniform}\{0, 1\}$
        \State $\calD_A, \calD_B, \calD_{\mathrm{BC}} \leftarrow
               \textsc{CollectRollouts}(\theta_A^{\mathrm{old}}, \theta_B^{\mathrm{old}}, \theta_E, r_{\mathrm{id}}, N_G)$
    \EndFor
    \State $\theta_A \leftarrow \theta_A - \eta\,\nabla_{\theta_A}\calL_A$
           \Comment{Entropy-regularised Alice PPO loss, \cref{eq:alice_loss}}
    \State $\theta_B \leftarrow \theta_B - \eta\,\nabla_{\theta_B}[\calL_{\mathrm{PPO}}^B + \beta\,\calL_{\mathrm{ABC}}]$
           \Comment{\cref{eq:bob_loss}}
    \State $\theta_E \leftarrow \theta_E - \eta\,\nabla_{\theta_E}\calL_{\mathrm{PPO}}^B$
           \Comment{Encoder updated via Bob's gradient}
\EndFor
\State \Return $\theta_B^*, \theta_E^*$ \Comment{Freeze these for Phase 2}
\end{algorithmic}
\end{algorithm}

\begin{algorithm}[H]
\caption{CollectRollouts (Phase 1)}
\label{alg:rollout_p1}
\begin{algorithmic}[1]
\Require $\theta_A^{\mathrm{old}}, \theta_B^{\mathrm{old}}, \theta_E$, $r_{\mathrm{id}}$, $N_G$
\State $\calD_A \leftarrow \emptyset$;\; $\calD_B \leftarrow \emptyset$;\;
       $\calD_{\mathrm{BC}} \leftarrow \emptyset$;\; $\xi \leftarrow \mathrm{True}$
\State $s_0 \leftarrow$ \textsc{Env.Reset()}
\For{goal $= 1, \ldots, N_G$}
    \State $\tau_A, \sstar \leftarrow$
           \textsc{GenerateAliceTraj}($\piA(\cdot;\theta_A^{\mathrm{old}})$, $s_0$, $r_{\mathrm{id}}$, $T_A$)
    \If{\textsc{GoalInvalid}($\sstar$)} \textbf{break} \EndIf
    \If{$\xi =$ True}
        \State $\tau_B, \xi \leftarrow$
               \textsc{GenerateBobTraj}($\piB'(\cdot;\theta_B^{\mathrm{old}}, \theta_E)$,
               $s_0$, $\sstar$, $r_{\mathrm{id}}$, $T_B$)
        \State $\calD_B \leftarrow \calD_B \cup \{\tau_B\}$
    \EndIf
    \State $r_A \leftarrow \textsc{ComputeAliceReward}(\xi, \sstar)$
    \State $\tau_A[-1].\mathrm{reward} \leftarrow r_A$;\;
           $\calD_A \leftarrow \calD_A \cup \{\tau_A\}$
    \If{$\xi = $ False} \Comment{Bob failed: Alice's trajectory is a useful demo}
        \State $\tau_{\mathrm{BC}} \leftarrow$ \textsc{RelabelDemonstration}($\tau_A$, $\sstar$,
               $\piB^{\mathrm{old}}$)
        \State $\calD_{\mathrm{BC}} \leftarrow \calD_{\mathrm{BC}} \cup \{\tau_{\mathrm{BC}}\}$
    \EndIf
    \State $s_0 \leftarrow$ last state of $\tau_B$ (or $\tau_A$ if Bob failed)
\EndFor
\State \Return $\calD_A, \calD_B, \calD_{\mathrm{BC}}$
\end{algorithmic}
\end{algorithm}

\subsection{Phase 2: Meta-ASP Manager Training}
\label{subsec:alg_phase2}

\begin{algorithm}[H]
\caption{Phase 2 — Manager Training via Meta-ASP}
\label{alg:phase2}
\begin{algorithmic}[1]
\Require $\theta_M, \theta_{MA}$ \Comment{Manager (meta-Bob) and meta-Alice parameters}
\Require $\theta_B^*, \theta_E^*$ \Comment{Frozen arm sub-policy and encoder from Phase 1}
\Require $\eta_M$, $\beta_M$, $\lambda$, $\mu$, $N_G^M$, $T_C$
\For{training step $= 1, 2, \ldots$}
    \State $\theta_M^{\mathrm{old}} \leftarrow \theta_M$;\;
           $\theta_{MA}^{\mathrm{old}} \leftarrow \theta_{MA}$
    \For{each parallel environment worker}
        \State $\calD_M, \calD_{MA}, \calD_{\mathrm{BC}}^M \leftarrow$
               \textsc{CollectManagerRollouts}($\theta_M^{\mathrm{old}}, \theta_{MA}^{\mathrm{old}},
               \theta_B^*, \theta_E^*, N_G^M, T_C$)
        \State $\calL_{\mathrm{col}}, \calL_{\mathrm{conf}} \leftarrow$
               \textsc{QuerySimPenalties}() \Comment{IsaacLab physics queries}
    \EndFor
    \State $\theta_{MA} \leftarrow \theta_{MA} - \eta_M\,\nabla_{\theta_{MA}}\calL_{MA}$
    \State $\theta_M \;\leftarrow \theta_M \;- \eta_M\,\nabla_{\theta_M}
           [\calL_{\mathrm{PPO}}^M + \beta_M\calL_{\mathrm{ABC}}^M
           + \lambda\calL_{\mathrm{col}} + \mu\calL_{\mathrm{conf}}]$
           \Comment{\cref{eq:manager_loss}}
\EndFor
\State \Return $\theta_M^*$
\end{algorithmic}
\end{algorithm}

% ══════════════════════════════════════════════════════════════════════════════
\section{IsaacLab Implementation Plan}
\label{sec:implementation}
% ══════════════════════════════════════════════════════════════════════════════

\subsection{Environment Setup}
\label{subsec:env_setup}

\begin{description}
    \item[Robot platform:] Two UR16e arms (or Franka Panda) with RobotIQ 2F-85 parallel
          grippers, mounted symmetrically facing a shared table surface. Arms are
          positioned so their reachable workspaces overlap in a central manipulation zone.

    \item[Scene:] 3--6 ShapeNet or YCB objects randomly placed on the table surface at
          episode start. Object types and counts are randomised for generalisation.

    \item[Simulation:] IsaacLab on PhysX, GPU-accelerated parallel environments
          ($N_{\mathrm{env}} = 4096$). Domain randomisation applied to object mass,
          friction, and initial placement.
\end{description}

\subsection{Action Space}
\label{subsec:action_space}

Following \citet{plappert2021asymmetric}, each arm uses a \textbf{6-dimensional} action
space:
\begin{itemize}[nosep]
    \item 3D relative gripper position (TCP Cartesian delta)
    \item 2D relative gripper rotation (wrist yaw + pitch delta)
    \item 1D gripper finger position (symmetric open/close)
\end{itemize}
Actions are discretised into 11 bins per dimension ($11^6$ discrete action space per arm)
and a multi-categorical distribution is learned. A PID controller converts TCP commands to
joint-space trajectories internally.

\subsection{Observation Space}
\label{subsec:obs_space}

Each arm sub-policy receives the observation vector in \cref{eq:obs}:

\begin{table}[h!]
\centering
\caption{Observation components per arm sub-policy.}
\begin{tabularx}{\textwidth}{llX}
\toprule
Component & Dim. & Description \\
\midrule
Own joint positions & 6 & Current joint angles \\
Own gripper position & 3 & TCP Cartesian position \\
Other arm joint positions & 6 & Second arm's joint angles (zero-padded in Phase 1) \\
Object state & $M \times 13$ & Per-object: position (3), rotation (4), velocity (3),
                               rot.\ velocity (3) \\
Goal embedding $\goalvec$ & $K$ & Manager-issued goal vector \\
Role identifier $r_{\mathrm{id}}$ & 1 & Binary flag: which joint block is ``me'' \\
\bottomrule
\end{tabularx}
\label{tab:obs}
\end{table}

The manager receives:

\begin{table}[h!]
\centering
\caption{Observation components for the manager.}
\begin{tabularx}{\textwidth}{llX}
\toprule
Component & Dim. & Description \\
\midrule
Arm 1 joint positions & 6 & \\
Arm 2 joint positions & 6 & \\
Arm 1 gripper position & 3 & TCP Cartesian \\
Arm 2 gripper position & 3 & TCP Cartesian \\
Object state & $M \times 13$ & Full scene object state \\
Task context $z_{\mathrm{task}}$ & $Z$ & Optional external task specification \\
\bottomrule
\end{tabularx}
\label{tab:obs_manager}
\end{table}

\subsection{Reward Functions}
\label{subsec:rewards}

\paragraph{Phase 1 (arm level).}
Bob receives sparse goal-conditioned reward from IsaacLab state queries:
\begin{equation}
    r_B^{\text{object}} = \mathbb{1}[D_{\mathrm{pos}}(o_k, g_k) < 0.04\,\mathrm{m}
    \;\wedge\; D_{\mathrm{rot}}(o_k, g_k) < 0.2\,\mathrm{rad}].
\end{equation}
Placing all $M$ objects correctly gives a bonus $r_B^{\text{success}} = 5$.
Moving a previously placed object away gives $r_B^{\text{disturb}} = -1$.
Alice receives $r_A = 5 + \mathbb{1}[\text{valid goal}]$ if Bob fails, 0 if Bob succeeds.

\paragraph{Phase 2 (manager level).}
Meta-Bob: sparse $\{0, 1\}$ binary reward for full workspace goal achievement.
Meta-Alice: $1 - R_{\text{meta-Bob}}$.
Both receive the physical penalties $\calL_{\mathrm{collision}}$ and
$\calL_{\mathrm{conflict}}$ as negative reward terms subtracted from their total reward.

\subsection{File Structure}
\label{subsec:files}

\begin{verbatim}
asyncDualPlayPPO/
├── tasks/
│   ├── dual_arm/
│   │   ├── env_cfg.py          # IsaacLab environment configuration
│   │   ├── obs_manager.py      # Observation builders (arm + manager)
│   │   ├── reward_phase1.py    # Bob/Alice sparse rewards
│   │   ├── reward_phase2.py    # Manager meta-ASP rewards + physics penalties
│   │   └── termination.py     # Episode termination conditions
├── algorithms/
│   ├── asp_phase1.py           # Phase 1 rollout + PPO + ABC
│   ├── asp_phase2.py           # Phase 2 rollout + manager PPO + ABC
│   └── goal_encoder.py         # Encoder E architecture
├── networks/
│   ├── arm_policy.py           # Shared sub-policy pi_B'
│   ├── manager_policy.py       # Manager pi_M (shared trunk + 2 heads)
│   └── alice_policy.py         # Alice pi_A and meta-Alice pi_A^M
├── utils/
│   ├── contact_sensor.py       # IsaacLab ContactSensor wrapper
│   ├── goal_validation.py      # meta-Alice goal validity checks
│   └── demo_filter.py          # ABC demonstration filtering
├── cfg/
│   ├── phase1_cfg.yaml         # Phase 1 hyperparameters
│   └── phase2_cfg.yaml         # Phase 2 hyperparameters
└── train.py                    # Entry point
\end{verbatim}

\subsection{Key Hyperparameters}
\label{subsec:hyperparams}

\begin{table}[h!]
\centering
\caption{Recommended hyperparameters based on \citet{plappert2021asymmetric,sukhbaatar2018learning}.}
\begin{tabular}{lll}
\toprule
Hyperparameter & Value & Source \\
\midrule
\multicolumn{3}{l}{\textit{Phase 1 (arm sub-policy)}} \\
Discount factor $\gamma$ & 0.998 & \cite{plappert2021asymmetric} \\
GAE $\lambda$ & 0.95 & \cite{plappert2021asymmetric} \\
PPO clip $\epsilon_{\mathrm{ppo}}$ & 0.2 & \cite{schulman2017proximal} \\
ABC clip $\epsilon_{\mathrm{abc}}$ & 0.2 & \cite{plappert2021asymmetric} \\
ABC weight $\beta$ & 0.5 & \cite{plappert2021asymmetric} \\
Alice entropy coeff.\ $\beta_A$ & 0.01 & \cite{sukhbaatar2018learning} \\
Alice steps $T_A$ & 100 & \cite{plappert2021asymmetric} \\
Bob max steps $T_B$ & 200 & \cite{plappert2021asymmetric} \\
Goals per episode $N_G$ & 5 & \cite{plappert2021asymmetric} \\
Goal embedding dim.\ $K$ & 8--16 & Proposed \\
Optimizer & Adam & \cite{kingma2014adam} \\
Learning rate $\eta$ & $3\times10^{-4}$ & \cite{plappert2021asymmetric} \\
Batch size & 4096 & \cite{plappert2021asymmetric} \\
\midrule
\multicolumn{3}{l}{\textit{Phase 2 (manager)}} \\
Manager goal period $T_C$ & $= T_A$ & \cite{sukhbaatar2018learning} \\
ABC weight $\beta_M$ & 0.5 & Proposed \\
Collision weight $\lambda_{\max}$ & 1.0 & Proposed \\
Collision weight $\lambda_{\min}$ & 0.1 & Proposed \\
Conflict weight $\mu$ & 0.5 & Proposed \\
Proximity threshold $d_{\mathrm{prox}}$ & 0.08\,m & Proposed \\
Manager goals $N_G^M$ & 5 & Following \cite{plappert2021asymmetric} \\
Past sampling (against old Alice) & 20\% & \cite{plappert2021asymmetric} \\
\bottomrule
\end{tabular}
\label{tab:hyperparams}
\end{table}

% ══════════════════════════════════════════════════════════════════════════════
\section{Evaluation Protocol}
\label{sec:eval}
% ══════════════════════════════════════════════════════════════════════════════

\subsection{Phase 1 Evaluation (Single-Arm)}
\label{subsec:eval_phase1}

Success rate on a fixed set of single-arm holdout tasks, following
\citet{plappert2021asymmetric}:
\begin{itemize}[nosep]
    \item \textbf{Push:} Move object(s) to target positions.
    \item \textbf{Flip:} Reorient object to target face.
    \item \textbf{Pick-and-place:} Lift one object to a target position in the air.
    \item \textbf{Stack:} Place one block atop another in the correct order.
\end{itemize}
Zero-shot success rate is measured by freezing Bob and evaluating on tasks never seen
during training (unseen object types, unseen goal configurations).

\subsection{Phase 2 Evaluation (Bimanual)}
\label{subsec:eval_phase2}

A set of bimanual-specific holdout tasks that \emph{require} both arms:
\begin{itemize}[nosep]
    \item \textbf{Handover:} Arm 1 picks up an object and transfers it to Arm 2's gripper.
    \item \textbf{Stabilise-and-insert:} Arm 1 holds a container; Arm 2 inserts an object.
    \item \textbf{Collaborative push:} Both arms push a large object together to a goal.
    \item \textbf{Bimanual stack:} Arms cooperate to place blocks in a tower configuration
          larger than either arm can achieve alone.
\end{itemize}

\paragraph{Metrics.}
\begin{itemize}[nosep]
    \item Zero-shot success rate (\%) per holdout task.
    \item Arm-arm collision count per episode.
    \item Goal proposal diversity: entropy of meta-Alice's workspace configuration distribution.
    \item Task curriculum progression: average difficulty of meta-Alice proposals over training.
\end{itemize}

% ══════════════════════════════════════════════════════════════════════════════
\section{Implementation Milestones}
\label{sec:milestones}
% ══════════════════════════════════════════════════════════════════════════════

\begin{table}[h!]
\centering
\caption{Implementation milestones with dependencies.}
\begin{tabularx}{\textwidth}{clXl}
\toprule
\# & Milestone & Deliverable & Depends on \\
\midrule
1 & Single-arm IsaacLab env & Environment with one arm, objects, obs/reward wrappers & --- \\
2 & Arm policy network & Sub-policy $\piB'$, encoder $\Enc$, Alice $\piA$, role-id obs & 1 \\
3 & Phase 1 ASP loop & Rollout collection, ABC buffer, PPO+ABC training & 2 \\
4 & Phase 1 evaluation & Holdout task success rates for single arm & 3 \\
5 & Dual-arm env & Two-arm scene, ContactSensor, shared workspace & 1 \\
6 & Manager network & Shared-trunk + 2-head manager, meta-Alice & 2, 5 \\
7 & Phase 2 ASP loop & Manager rollouts, sim penalties, meta-ABC & 4, 6 \\
8 & Phase 2 evaluation & Bimanual holdout tasks & 7 \\
9 & Ablation studies & Disable ABC, collision penalty, role-id & 8 \\
\bottomrule
\end{tabularx}
\label{tab:milestones}
\end{table}

% ══════════════════════════════════════════════════════════════════════════════
\section{Known Challenges and Open Questions}
\label{sec:challenges}
% ══════════════════════════════════════════════════════════════════════════════

\begin{enumerate}
    \item \textbf{Distribution shift at deployment.} Sub-policies see zero-padded
          other-arm observations during Phase~1 but non-zero values at deployment.
          Mitigation: gradually replace zeros with recorded arm trajectories from earlier
          Phase~1 episodes during the final training stage.

    \item \textbf{Role specialisation from a shared policy.} A single policy with role-id
          may converge to symmetric behaviour. Monitoring whether the manager learns to
          send consistently different $\goalvec^{(1)}$ vs.\ $\goalvec^{(2)}$ is critical.
          If symmetry is observed, an asymmetric initialisation or role-specific fine-tuning
          stage may be required.

    \item \textbf{Manager curriculum coverage.} meta-Alice must discover bimanual goals
          that genuinely require both arms. Alice's entropy regularisation
          \cref{eq:alice_loss} may be insufficient if single-arm-solvable goals dominate
          the workspace configuration space. A coverage metric (fraction of objects touched
          by each arm during meta-Alice's turn) can be used as an auxiliary reward.

    \item \textbf{T\textsubscript{C} mismatch.} If one arm's sub-goal is achieved in
          significantly fewer than $T_C$ steps, it continues executing with a stale goal.
          Mitigation: the sub-policy can be trained with a ``hold position'' goal type that
          it automatically invokes after success detection.

    \item \textbf{Meta-ASP resettability.} The manager-level ASP requires that
          meta-Alice's trajectory reaches $\sstar_{\mathrm{ws}}$ from $s_0$ using the
          same frozen arms as meta-Bob. If meta-Alice learns to exploit physics in ways
          the arms cannot replicate deterministically (e.g., sliding objects with gripper
          forces that are highly sensitive to friction randomisation), the ABC demonstrations
          become misleading. Demonstration filtering (retain only trajectories that meta-Bob
          fails) partially mitigates this.
\end{enumerate}

% ══════════════════════════════════════════════════════════════════════════════
\section{Ablation Study Design}
\label{sec:ablations}
% ══════════════════════════════════════════════════════════════════════════════

Following the ablation structure of \citet{plappert2021asymmetric}:

\begin{table}[h!]
\centering
\caption{Planned ablation studies.}
\begin{tabular}{ll}
\toprule
Ablation & Component disabled \\
\midrule
No Phase 1 pre-training & Manager trained from scratch without frozen sub-policies \\
No ABC at arm level & $\beta = 0$ in \cref{eq:bob_loss} \\
No ABC at manager level & $\beta_M = 0$ in \cref{eq:manager_loss} \\
No collision penalty & $\lambda = 0$ in \cref{eq:manager_loss} \\
No role identifier & $r_{\mathrm{id}}$ removed from observation \cref{eq:obs} \\
Single-goal meta-ASP & $N_G^M = 1$ instead of 5 \\
External reward manager & Meta-ASP replaced by task reward supervision \\
\bottomrule
\end{tabular}
\label{tab:ablations}
\end{table}

% ══════════════════════════════════════════════════════════════════════════════
\bibliographystyle{plainnat}
\begin{thebibliography}{99}

\bibitem[Andrychowicz et~al.(2017)]{andrychowicz2017hindsight}
Andrychowicz, M., Wolski, F., Ray, A., Schneider, J., Fong, R., Welinder, P.,
McGrew, B., Tobin, J., Abbeel, P., and Zaremba, W.
\newblock Hindsight experience replay.
\newblock In \emph{Advances in Neural Information Processing Systems}, 2017.

\bibitem[IsaacLab(2023)]{isaaclab2023}
Isaac-Sim Team.
\newblock IsaacLab: A Framework for Robot Learning Research.
\newblock \url{https://isaac-sim.github.io/IsaacLab/main/index.html}, 2023.

\bibitem[Kaelbling(1993)]{kaelbling1993learning}
Kaelbling, L.~P.
\newblock Learning to achieve goals.
\newblock In \emph{Proceedings of the 13th International Joint Conference on
  Artificial Intelligence}, 1993.

\bibitem[Kingma and Ba(2014)]{kingma2014adam}
Kingma, D.~P. and Ba, J.
\newblock Adam: A method for stochastic optimization.
\newblock \emph{arXiv preprint arXiv:1412.6980}, 2014.

\bibitem[Nachum et~al.(2018)]{nachum2018data}
Nachum, O., Gu, S., Lee, H., and Levine, S.
\newblock Data-efficient hierarchical reinforcement learning.
\newblock In \emph{Advances in Neural Information Processing Systems}, 2018.

\bibitem[Plappert et~al.(2021)]{plappert2021asymmetric}
Plappert, M., Sampedro, R., Xu, T., Akkaya, I., Kosaraju, V., Welinder, P.,
D'Sa, R., Petron, A., Pinto, H.~P.~de~O., Paino, A., Noh, H., Weng, L., Yuan, Q.,
Chu, C., and Zaremba, W.
\newblock Asymmetric self-play for automatic goal discovery in robotic manipulation.
\newblock \emph{arXiv preprint arXiv:2101.04882}, 2021.

\bibitem[Schulman et~al.(2017)]{schulman2017proximal}
Schulman, J., Wolski, F., Dhariwal, P., Radford, A., and Klimov, O.
\newblock Proximal policy optimization algorithms.
\newblock \emph{arXiv preprint arXiv:1707.06347}, 2017.

\bibitem[Sukhbaatar et~al.(2018a)]{sukhbaatar2018learning}
Sukhbaatar, S., Denton, E., Szlam, A., and Fergus, R.
\newblock Learning goal embeddings via self-play for hierarchical reinforcement
  learning.
\newblock \emph{arXiv preprint arXiv:1811.09083}, 2018.

\bibitem[Sukhbaatar et~al.(2018b)]{sukhbaatar2018intrinsic}
Sukhbaatar, S., Lin, Z., Kostrikov, I., Synnaeve, G., Szlam, A., and Fergus, R.
\newblock Intrinsic motivation and automatic curricula via asymmetric self-play.
\newblock In \emph{International Conference on Learning Representations}, 2018.

\bibitem[Vezhnevets et~al.(2017)]{vezhnevets2017feudal}
Vezhnevets, A.~S., Osindero, S., Schaul, T., Heess, N., Jaderberg, M.,
Silver, D., and Kavukcuoglu, K.
\newblock Feudal networks for hierarchical reinforcement learning.
\newblock In \emph{International Conference on Machine Learning}, 2017.

\end{thebibliography}

\end{document}
